\SourceLeaf{335}

If values of $\mathfrak{B}, \mathfrak{C}$ are prescribed that fail to satisfy these conditions or any one of them, the problem entails a physical impossibility.

Conversely, if the prescribed values satisfy the conditions, the problem is possible. That is, the two auxiliary circles meet at a point that lies in the suitable segment of each circle. In case 1), if the signs of those sines are $+\;+\;+$, their intersection cannot lie outside the main circle, where the first two sines would have to be negative and the third positive. The first two conditions would then require $\mathfrak{B}+180^0$, $\mathfrak{C}+180^0$ to take the place of $\mathfrak{B}, \mathfrak{C}$, which in turn contradicts the third condition. Thus the intersection lies inside the main circle, where the suitable segments lie.

Similarly, 2) for the signs $-\;-\;+$ one proves that the intersection can only lie outside, where the suitable segments then lie.

\begin{center}
\begin{tabular}{c@{\hspace{2em}}cc}
If the signs are & \multicolumn{2}{c}{the intersection lies} \\
 & in the suitable segment & in the unsuitable one \\
$+\;+\;+$ & circle on $a$ & [circle on] $b, c$ \\
$+\;-\;+$ & $b$ & $a, c$ \\
$+\;-\;-$ & $c$ & $a, b$ \\
$-\;-\;-$ & $a$ & $b, c$ \\
$-\;+\;+$ & $c$ & $a, b$ \\
$-\;+\;-$ & $b$ & $a, c$ \\
\end{tabular}
\end{center}

The elementary construction for \textsc{Pothenot}'s problem consists in determining points on the straight line through (0) perpendicular to 10. For one such point the polar coordinates relative to (0) are
\[ \text{Distance } \frac{(12)}{\sin(02-01)}, \qquad \text{Direction } 12+90^0-(02-01). \]

Another new method of construction starts by seeking a point other than (0), namely a point $(*)$ that lies on the straight line 10 at the distance $(1*)=\frac{1}{(10)}$. This point is determined by the intersections of straight lines, with one fewer line than there are objects sighted from (0). One of these
